\documentclass[journal]{vgtc}

\onlineid{0}
\vgtccategory{Research}

\title{Pseudocode-Map Substitution Grows with Expertise}

\author{%
  Razieh Fathi and Ian Handy
}

\authorfooter{%
  \item Razieh Fathi is with [CONFIRM: institution]. E-mail: [CONFIRM].
  \item Ian Handy is with [CONFIRM: institution]. E-mail: [CONFIRM].
}

\abstract{%
  This study examined how novice and experienced programmers attended
  to Pseudocode and maps presented simultaneously while viewing video
  tutorials about breadth-first search (BFS) and depth-first search
  (DFS) graph traversals. Using eye tracking technology, we recorded
  where each of our 62 participants directed their gaze toward the
  Pseudocode or the Map during each of the 117 trials they completed.
  As programming knowledge increased from novice to intermediate to
  advanced, the correlations between how much time was spent gazing
  at the Pseudocode versus the Map became significantly more negative.
  This represents a type of Expertise Reversal Effect. While there
  was a similar trend in terms of how many times each area of
  interest was viewed, the trend did not reach statistical
  significance. Therefore, although these findings support the idea
  that a design intended to be optimal for novice learners might be
  suboptimal for expert learners, they also suggest caution regarding
  generalizing results based upon findings obtained from a limited
  number of novice learners.
}

\keywords{eye-tracking, algorithm visualization, expertise reversal effect, dual representation, representation substitution}

\graphicspath{{figs/}{figures/}{pictures/}{images/}{./}}

\usepackage{amsmath}
\usepackage{booktabs}
\usepackage{mathptmx}
\usepackage{subcaption}

\begin{document}

\firstsection{Introduction}

\maketitle

The ability to combine multiple representations has been shown to
enhance learning~\cite{Mayer2001Multimedia}. However, the extent to
which such combinations are effective will depend upon the prior
experiences of those who receive instruction. Both Cognitive Load
Theory (CLT)~\cite{Sweller1988Cognitive} and the Expertise Reversal
Effect (ERE)~\cite{Kalyuga2003Expertise} propose that instructional
aids, which are beneficial for novices, do not provide additional
benefits for experienced learners. CLT states that the primary source
of cognitive load is due to the complexity of the content being
learned. When an experienced learner receives instruction via an aid
designed for novices, ERE suggests that the added load created by the
aid will outweigh its potential benefits. ERE proposes that as
learners gain experience and form a schema related to the subject
matter being taught, aids that previously supported the learning
process become unnecessary.

Therefore, in order to determine whether such a trade-off occurs and
whether the trade-off varies depending upon the level of expertise of
a given learner, we calculated the Spearman correlation coefficient
between the total amount of time spent gazing at the Pseudocode and
the total amount of time spent gazing at the Map. These calculations
allowed us to assess the degree to which learners substituted one
representation for another.

Contributions: (1)~we found significant evidence that the amount of
time experienced learners spend gazing at Pseudocode and/or a Map
decreases as they begin to substitute one for the other; (2)~we
provided a correlation-based method for quantifying representation
substitution; and (3)~we suggested implications for designing
split-panel visualizations that cater specifically to certain levels
of expertise.


\section{Literature Review}

CTML (Cognitive Theory of Multimedia Learning), developed by
Mayer~\cite{Mayer2001Multimedia}, posits that presenting verbal and
pictorial information via separate pathways into working memory
allows for better integration than when both types of information are
combined. Visualizations involving split panels apply this concept
directly to the presentation of algorithms. CLT provides a
theoretical basis for understanding why adding a second representation
to an existing representation would increase cognitive load.
Specifically, CLT states that when the intrinsic difficulty of a task
is already high, a secondary representation will result in greater
cognitive load than a single
representation~\cite{Sweller1988Cognitive,Sweller1994Cognitive}.
Finally, ERE expands on CLT by proposing that instructional aids that
assist novices become redundant once learners have formed a schema
related to the subject matter being
studied~\cite{Kalyuga2003Expertise}.

A meta-analysis conducted by
Gegenfurtner~\cite{Gegenfurtner2011Expertise} revealed that
intermediate learners tend to focus more than both novice and expert
learners, which challenges the traditional view that as prior
knowledge increases so too does focused attention.
Bednarik~\cite{Bednarik2012Expertise} demonstrated that experts and
novices differ in their attentional strategies when debugging
algorithms. Experts verify their mental models whereas novices
oscillate back-and-forth between panels without relying upon a
strategic approach to problem solving. Similarly, Richter and
Scheiter~\cite{Richter2019Studying} discovered that
low-prior-knowledge learners use signals to coordinate text and
picture areas-of-interest (AOI) when providing cues; however,
high-prior-knowledge learners do not benefit from cues.

Neither Bednarik nor Richter and Scheiter tested three levels of
expertise. Furthermore, neither study measured correlation between
dwell times as a metric. To demonstrate a substitution gradient
requires both measures.


\section{Methods}

\subsection{Using METAL as an Interactive Algorithm Visualization}

Transfer of learning occurs when learners apply previously acquired
knowledge to new problems or contexts~\cite{Olson2015Transfer}. Moving
beyond surface-level schemas often requires additional instructional
support, so that learners can attend to the reasoning behind each
step of a solution rather than its form alone. Interactive
visualizations embedded in example-based hypermedia have been
suggested as one possible support~\cite{VanMerrienboer2003Taking}.

We use the Map-based Educational Tools for Algorithmic Learning
(METAL) project~\cite{Teresco2018METAL} as our visualization
platform. METAL renders algorithm execution on Leaflet-based
geographic maps alongside color-coded data-structure tables, and has
been reported to increase student engagement in coursework that uses
it~\cite{Teresco2018METAL}. For the present study we use METAL's
visualizations for breadth-first search (BFS) and depth-first search
(DFS) applied to a single map dataset---the Delaware road network
(henceforth DE; 735 vertices, 879 edges; August~2022 TMG snapshot).
Figure~\ref{fig:metal} shows a representative frame from one trial.

\begin{figure}[tb]
  \centering
  \includegraphics[width=\linewidth, alt={Screenshot of the METAL HDX interface running a breadth-first search traversal on the Delaware road network, with a group-aggregated gaze-density heatmap overlaid on the map panel.}]{fig_metal_screenshot}
  \caption{%
    The METAL HDX interface running a BFS traversal on the Delaware
    (DE) road network. The left panels show the
    algorithm-visualization status (including the BFS discovered queue
    and partial spanning tree) and the per-edge traversal table; the
    right panel shows the geospatial map. Warm colors overlaid on the
    map denote aggregated gaze density across Group~1 participants
    (darker red = more fixations).
  }
  \label{fig:metal}
\end{figure}

\subsection{Design of the Study}

The study employed a two-factor between/within design. The
between-subject factor was prior programming experience (three
groups: novice, intermediate, advanced; see~\Cref{sec:subjects}). The
within-subject factor was algorithm (BFS, DFS); every participant
viewed both algorithms in the same session. Both BFS and DFS are
standardly covered in introductory data-structures curricula, and are
discoverable, if not always comprehended, by computer-science majors
and non-majors alike. Trial order and other counterbalancing of BFS
vs.\ DFS presentation are [CONFIRM].

\subsection{Subjects}
\label{sec:subjects}

Sixty-two participants were recruited at [CONFIRM: institution]; the
study was approved under IRB protocol [CONFIRM: number]. Gaze was
recorded with a Tobii eye tracker sampled at 60~Hz as participants
viewed animated BFS and DFS traversals in the METAL AV system.

We assigned each participant to one of three prior-knowledge groups,
according to [CONFIRM: measurement method]:
\begin{itemize}
  \item \textbf{G1 --- novice:} no background in computer science.
  \item \textbf{G2 --- intermediate:} brief exposure to computer
    science (e.g., preceding module of coursework).
  \item \textbf{G3 --- advanced:} multiple years of formal
    computer-science coursework.
\end{itemize}

Every participant viewed both the BFS and DFS animations. With
exclusions for quality control, the analyzable set comprised 117
trials ($n$=59 BFS, $n$=58 DFS), with counts per group for analyzable
set G1=19, G2=21, G3=20. All the results reported below are
conditioned on the DE stimulus rendered in the METAL style; the IRN
stimulus present in the aggregate dataset is not included in this
analysis.

\subsection{Stimuli and AOIs}
\label{sec:stimuli}

Each trial consisted of a split-screen METAL stimulus: a Pseudocode
panel on the left, showing the algorithm's current step, and a
geospatial Map on the right, showing what nodes and edges light up as
the algorithm traverses through them. Two AOIs were defined on the DE
stimulus---Pseudocode and Map. All participants viewed the BFS
animation (\texttt{DE-bft.wmv}) and the DFS animation
(\texttt{DE-dft.wmv}); length of each trial was [CONFIRM: number].

For each AOI we made two primary measurements---Fixation Count~(FC),
and Total Fixation Duration~(TFD) across all fixations in that AOI.
We derived from these the ratio
$\text{TFD}_\text{pseudo}/\text{TFD}_\text{map}$ of TFD for Pseudocode
to TFD for Map, scanner index
$\text{VC}_\text{pseudo}/\text{TFD}_\text{pseudo}$, average depth
$\text{TFD}_\text{pseudo}/\text{FC}_\text{pseudo}$, and switching
rate
$((\text{VC}_\text{pseudo}+\text{VC}_\text{map})/
(\text{TFD}_\text{pseudo}+\text{TFD}_\text{map}))$.

\subsection{Procedure}

Participants completed the study individually in a single session
conducted at [CONFIRM: location] between October and November 2022.
After signing an informed consent form, each participant was seated
in front of the eye-tracking display and completed a [CONFIRM:
calibration procedure]. Then, participants viewed both their BFS and
DFS animations [CONFIRM: order]---each shown once, per participant.
Either as part of, or following this task, we also [CONFIRM: if
additional practice trials, comprehension checks, or other post-task
questionnaires were administered]. Sessions took approximately
[CONFIRM: number] minutes per participant. [CONFIRM: compensation].
We recorded data using [CONFIRM: Tobii software and version] and
exported data per-AOI aggregated, for offline analysis.


\section{Results}

\subsection{A Non-Monotonic Pattern of Engagement in Pseudocode}
\label{sec:inverted}

G2 viewed the pseudocode panel longer than G1 and G3. The median
pseudocode-to-map ratio was 12.33 in G2. G1 and G3 showed medians of
9.57 and 6.12 respectively. This represents an inverse U relationship
with respect to both Total Fixation Duration~(TFD) and average
fixation depth. Only G1 was different ($p$=.040) when comparing BFS
and DFS. However, there was no difference ($p$>.50) when comparing
G2 and G3. Thus, prior knowledge appears to override the impact of
visual organization in determining where participants allocate their
gaze. The same shape persists when conditioning on the DFS stimulus
alone, with G2 highest and G3 lowest (\Cref{fig:results}(a)).

\subsection{Monotonic Substitution with Respect to Expertise}
\label{sec:reversal}

As expertise increased, the Spearman rank order correlation
coefficient ($\rho$) between pseudocode view TFD and map view TFD
became more negative. $\rho$(G1)=$-$.034, NS, $n$=19;
$\rho$(G2)=$-$.457, $p$<.050, $n$=21; $\rho$(G3)=$-$.631, $p$=.003,
$n$=20. Fisher z supported the end point comparison (G1 vs G3,
$p$<.050); however, neither of the middle comparisons reached
statistical significance (G1 vs G2 $p$=.180; G2 vs G3 $p$=.470).

Similarly, the rank order correlation between pseudocode view
fixation count and map view fixation count followed the same trend
but had lower values. $\rho$(G1)=+.042, NS; $\rho$(G2)=+.326, NS;
$\rho$(G3)=$-$.161, NS. All three correlations failed to reach
statistical significance ($p$>.100). Similarly, pairwise Fisher z on
FC failed to reach significance ($p$>.10): G2 vs G3 ($p$=.130); G1
vs G3 ($p$=.550); G1 vs G2 ($p$=.380). Partial Spearman controlling
for group and algorithm yielded directional relationships similar to
those found for TFD. \Cref{fig:results} shows both the inverted-U
and the lead TFD scatter.

\begin{figure*}[tb]
  \centering
  \begin{subfigure}[t]{0.40\linewidth}
    \includegraphics[width=\linewidth, alt={Boxplot of pseudocode-to-map TFD ratio by expertise group on the DFS Metal stimulus.}]{fig_ratio_inverted_u_dfs}
    \caption{Inverted-U in pseudocode/map TFD ratio: G2 highest, G3 lowest.}
  \end{subfigure}\hfill
  \begin{subfigure}[t]{0.58\linewidth}
    \includegraphics[width=\linewidth, alt={Scatter plot of TFD pseudocode vs map for G1 G2 G3 on DFS Metal.}]{fig_tfd_substitution_dfs}
    \caption{Primary result: TFD substitution. G3 $\rho$=$-$.66, $p$=.002; G1 flat.}
  \end{subfigure}
  \caption{DFS DE Metal results. (a) Per-participant pseudocode/map TFD
    ratio shows the inverted-U with G2 at the apex.
    (b) Per-participant TFD on pseudocode vs.\ TFD on map: as
    expertise rises the regression slope rotates from flat (G1) to
    strongly negative (G3, $\rho$=$-$.66, $p$=.002).}
  \label{fig:results}
\end{figure*}

The TFD findings match the predictions of ERE regarding visual search
behavior. Novices use both panels as separate units of information
processing space; intermediates begin to substitute one unit for
another; and experts substitute most heavily. As a result, time
devoted to viewing one unit becomes time that is subtracted from
viewing another unit. The predicted coordination effect of CTML due
to the dual channel nature of the layout is absent in every group's
dwell time measures.

\subsection{Attentional Coupling Increases with Increasing Expertise}

Attentional coupling increases with increasing expertise. The mean
absolute value of Spearman's rank order correlation coefficient
($\rho$) across all possible metric pairs: G1=.294; G2=.330; G3=.369.
Participants without prior experience with visual representations of
algorithms demonstrated the least coupled metric pairs. For example,
the ratio to scanner index pair ranged from $r$=$-$.04 in G1 to
$r$=$-$.67 in G3. Restricted to the DFS Metal corpus the gradient
is sharper still (mean $|\rho|$ = .368, .483, .605 for G1, G2, G3).
Scanner index, fixation depth, and switching rate did not differ by
group (all $p$>.20), ruling out a global pacing explanation.


\section{Discussion}

\subsection{Mechanisms}

The monotonic increase in the magnitude of negative TFD correlation
across the three groups is consistent with CLT and ERE: novices
without a schema (G1) consume working memory learning both the
algorithm and the visualization, so dwell on one panel is unrelated
to dwell on the other; intermediates (G2) develop predictable per-panel
reading patterns, so dwell begins to trade off; advanced learners (G3)
treat each panel as a sufficient cue and substitute one for the other.
The inverted-U in dwell ratio (\Cref{sec:inverted}) and the monotonic
substitution gradient (\Cref{sec:reversal}) thus describe the same
process at the population and per-participant scales respectively.

\subsection{Design Implications}

Since substitution varies with respect to expertise level, each
group of learners uses the split-panel design differently than
others. G1: Both panels are used independently. Scaffolding (e.g.,
line highlighting, annotation, staged display) may be necessary to
facilitate coordination between panels.

G2: Panels are partially substitutable with respect to each other
while still being engaged simultaneously. This profile fits well
with the split-panel format.

G3: The pseudocode panel serves as a verification check for either
representation alone.

Metrics related to attention aggregated across all levels of
expertise provide average values that characterize none of the
groups accurately.

\subsection{Limitations}

It is expected that roughly five of the 105 Fisher z tests would
reach $p$<.05 by chance alone. While the main TFD contrast (G1 vs
G3, $p$=.04) survives that threshold with some modest correction
applied, neither of the FC sign reversals reached statistical
significance at the participant level. Learner expertise groups were
not randomly assigned, thus group-level differences may be reflective
of either general cognitive ability or learner motivation. All
visualizations utilized Metal style. Group sizes of 19 to 21 limited
power for detecting small effects.

\subsection{Conclusion}

Time allocated to viewing pseudocode versus maps trades off
positively with programming expertise level. Time allocations are
essentially independent for novice programmers; intermediate and
advanced programmers substitute one for the other increasingly with
greater expertise level. Layouts validated at one level of expertise
may not generalize; aggregating attentional metrics across levels of
expertise hides this substitution gradient.


\bibliographystyle{abbrv-doi-hyperref}
\bibliography{paper}


\end{document}
